\documentclass[12pt]{ximera}
\typeout{Start loading xmPreamble.tex}%

% Add here extra macro's that are loaded automatically by all documents of claas 'ximera' or 'xourse' in this repo

%%
%%  Example:
%%
% \newcommand{\R}{\mathbb{R}}
\usepackage{tikz}
\usetikzlibrary{positioning, arrows.meta}
\usepackage{multicol}
\usepackage{enumitem}
\usepackage{caption}
\usepackage{amsmath}
\usepackage{amsthm}
\usepackage{fontenc}

%% ---------------------------------------------------------------
%% Shared worksheet macros.
%%
%% These came from the Summer 2026 worksheet set, where each worksheet carried its
%% own copy. They have to live here instead: a xourse inlines every activity into a
%% single document, so duplicate \newcommand definitions across activities collide,
%% and the xourse gobbles \input so a per-topic macro file never loads.
%%
%% They use \providecommand, NOT \newcommand. ximera.cls loads this file automatically,
%% and every activity also carries an explicit \typeout{Start loading xmPreamble.tex}%

% Add here extra macro's that are loaded automatically by all documents of claas 'ximera' or 'xourse' in this repo

%%
%%  Example:
%%
% \newcommand{\R}{\mathbb{R}}
\usepackage{tikz}
\usetikzlibrary{positioning, arrows.meta}
\usepackage{multicol}
\usepackage{enumitem}
\usepackage{caption}
\usepackage{amsmath}
\usepackage{amsthm}
\usepackage{fontenc}

%% ---------------------------------------------------------------
%% Shared worksheet macros.
%%
%% These came from the Summer 2026 worksheet set, where each worksheet carried its
%% own copy. They have to live here instead: a xourse inlines every activity into a
%% single document, so duplicate \newcommand definitions across activities collide,
%% and the xourse gobbles \input so a per-topic macro file never loads.
%%
%% They use \providecommand, NOT \newcommand. ximera.cls loads this file automatically,
%% and every activity also carries an explicit \typeout{Start loading xmPreamble.tex}%

% Add here extra macro's that are loaded automatically by all documents of claas 'ximera' or 'xourse' in this repo

%%
%%  Example:
%%
% \newcommand{\R}{\mathbb{R}}
\usepackage{tikz}
\usetikzlibrary{positioning, arrows.meta}
\usepackage{multicol}
\usepackage{enumitem}
\usepackage{caption}
\usepackage{amsmath}
\usepackage{amsthm}
\usepackage{fontenc}

%% ---------------------------------------------------------------
%% Shared worksheet macros.
%%
%% These came from the Summer 2026 worksheet set, where each worksheet carried its
%% own copy. They have to live here instead: a xourse inlines every activity into a
%% single document, so duplicate \newcommand definitions across activities collide,
%% and the xourse gobbles \input so a per-topic macro file never loads.
%%
%% They use \providecommand, NOT \newcommand. ximera.cls loads this file automatically,
%% and every activity also carries an explicit \input{../xmPreamble}, so this file is
%% read twice. That was harmless while it held only \usepackage lines (LaTeX skips a
%% package it has already loaded) but a second \newcommand is a hard error.
%%
%%   \blank[width]        fill-in-the-blank rule (default 2.5cm)
%%   \showwork            "show and explain your work" reminder
%%   \showworkline        ... plus which schedule line was used
%%   \showworkyear        ... plus which year's schedule was used
%%   \schedhead{...}      centred bold caption above a tiered-rate table
%%   \samesquares[scale]{left}{right}   two equal-area squares, labelled
%%   \samecubes[scale]{left}{right}     two equal-volume cubes, labelled
%% ---------------------------------------------------------------

\providecommand{\blank}[1][2.5cm]{\underline{\hspace{#1}}}
\providecommand{\showwork}{\textbf{REMEMBER TO SHOW AND EXPLAIN YOUR WORK.}}
\providecommand{\showworkline}{\textbf{REMEMBER TO SHOW AND EXPLAIN YOUR WORK.
  Explanation should include how and why you know which line of the
  schedule was used to calculate this amount.}}
\providecommand{\showworkyear}{\begin{sloppypar}\textbf{REMEMBER TO SHOW AND
  EXPLAIN YOUR WORK. Explanation should include how and why you know which
  line of the year's tax schedule was used to calculate this tax
  amount.}\end{sloppypar}}
\providecommand{\schedhead}[1]{\begin{center}\textbf{#1}\end{center}}

\providecommand{\samesquares}[3][1]{%
\begin{center}
{\footnotesize These squares are the \textbf{same size}.}

\vspace{1.5mm}
\begin{tikzpicture}[scale=#1,every node/.style={scale=#1}]
  \draw (0,0) rectangle (2.4,2.4);
  \node[anchor=north west] at (0.15,2.25) {Area:};
  \node[anchor=east]  at (-0.15,1.2) {#2};
  \node[anchor=north] at (1.2,-0.15) {#2};
  \begin{scope}[xshift=4.6cm]
    \draw (0,0) rectangle (2.4,2.4);
    \node[anchor=north west] at (0.15,2.25) {Area:};
    \node[anchor=east]  at (-0.15,1.2) {#3};
    \node[anchor=north] at (1.2,-0.15) {#3};
  \end{scope}
\end{tikzpicture}
\end{center}}

\providecommand{\samecubes}[3][1]{%
\begin{center}
{\footnotesize These cubes are the \textbf{same size}.}

\vspace{1.5mm}
\begin{tikzpicture}[scale=#1,every node/.style={scale=#1}]
  \foreach \dx in {0,5.2} {
    \begin{scope}[xshift=\dx cm]
      \draw (0,0) rectangle (2.0,2.0);
      \draw (0,2.0) -- (0.65,2.6) -- (2.65,2.6) -- (2.0,2.0);
      \fill[black!12] (2.0,0) -- (2.65,0.6) -- (2.65,2.6) -- (2.0,2.0) -- cycle;
      \draw (2.0,0) -- (2.65,0.6) -- (2.65,2.6) -- (2.0,2.0);
      \node[anchor=north west] at (0.1,1.9) {Volume:};
    \end{scope}
  }
  \node[anchor=east]  at (-0.15,1.0) {#2};
  \node[anchor=north] at (1.0,-0.15) {#2};
  \node[anchor=west]  at (2.25,0.4) {#2};
  \node[anchor=east]  at (5.05,1.0) {#3};
  \node[anchor=north] at (6.2,-0.15) {#3};
  \node[anchor=west]  at (7.45,0.4) {#3};
\end{tikzpicture}
\end{center}}

%% NOTE: do NOT add \usepackage to individual activity .tex files.
%% When a xourse inlines an activity it gobbles \documentclass and \input, but a bare
%% \usepackage still executes -- after \begin{document} -- and errors out.
%% All shared packages and macros belong here.
%%
%% ximera.cls already loads: amsmath, amssymb, amsthm, cancel, enumitem, graphicx,
%% hyperref, listings, pgfplots, tikz, url, xcolor. No need to re-load those.
, so this file is
%% read twice. That was harmless while it held only \usepackage lines (LaTeX skips a
%% package it has already loaded) but a second \newcommand is a hard error.
%%
%%   \blank[width]        fill-in-the-blank rule (default 2.5cm)
%%   \showwork            "show and explain your work" reminder
%%   \showworkline        ... plus which schedule line was used
%%   \showworkyear        ... plus which year's schedule was used
%%   \schedhead{...}      centred bold caption above a tiered-rate table
%%   \samesquares[scale]{left}{right}   two equal-area squares, labelled
%%   \samecubes[scale]{left}{right}     two equal-volume cubes, labelled
%% ---------------------------------------------------------------

\providecommand{\blank}[1][2.5cm]{\underline{\hspace{#1}}}
\providecommand{\showwork}{\textbf{REMEMBER TO SHOW AND EXPLAIN YOUR WORK.}}
\providecommand{\showworkline}{\textbf{REMEMBER TO SHOW AND EXPLAIN YOUR WORK.
  Explanation should include how and why you know which line of the
  schedule was used to calculate this amount.}}
\providecommand{\showworkyear}{\begin{sloppypar}\textbf{REMEMBER TO SHOW AND
  EXPLAIN YOUR WORK. Explanation should include how and why you know which
  line of the year's tax schedule was used to calculate this tax
  amount.}\end{sloppypar}}
\providecommand{\schedhead}[1]{\begin{center}\textbf{#1}\end{center}}

\providecommand{\samesquares}[3][1]{%
\begin{center}
{\footnotesize These squares are the \textbf{same size}.}

\vspace{1.5mm}
\begin{tikzpicture}[scale=#1,every node/.style={scale=#1}]
  \draw (0,0) rectangle (2.4,2.4);
  \node[anchor=north west] at (0.15,2.25) {Area:};
  \node[anchor=east]  at (-0.15,1.2) {#2};
  \node[anchor=north] at (1.2,-0.15) {#2};
  \begin{scope}[xshift=4.6cm]
    \draw (0,0) rectangle (2.4,2.4);
    \node[anchor=north west] at (0.15,2.25) {Area:};
    \node[anchor=east]  at (-0.15,1.2) {#3};
    \node[anchor=north] at (1.2,-0.15) {#3};
  \end{scope}
\end{tikzpicture}
\end{center}}

\providecommand{\samecubes}[3][1]{%
\begin{center}
{\footnotesize These cubes are the \textbf{same size}.}

\vspace{1.5mm}
\begin{tikzpicture}[scale=#1,every node/.style={scale=#1}]
  \foreach \dx in {0,5.2} {
    \begin{scope}[xshift=\dx cm]
      \draw (0,0) rectangle (2.0,2.0);
      \draw (0,2.0) -- (0.65,2.6) -- (2.65,2.6) -- (2.0,2.0);
      \fill[black!12] (2.0,0) -- (2.65,0.6) -- (2.65,2.6) -- (2.0,2.0) -- cycle;
      \draw (2.0,0) -- (2.65,0.6) -- (2.65,2.6) -- (2.0,2.0);
      \node[anchor=north west] at (0.1,1.9) {Volume:};
    \end{scope}
  }
  \node[anchor=east]  at (-0.15,1.0) {#2};
  \node[anchor=north] at (1.0,-0.15) {#2};
  \node[anchor=west]  at (2.25,0.4) {#2};
  \node[anchor=east]  at (5.05,1.0) {#3};
  \node[anchor=north] at (6.2,-0.15) {#3};
  \node[anchor=west]  at (7.45,0.4) {#3};
\end{tikzpicture}
\end{center}}

%% NOTE: do NOT add \usepackage to individual activity .tex files.
%% When a xourse inlines an activity it gobbles \documentclass and \input, but a bare
%% \usepackage still executes -- after \begin{document} -- and errors out.
%% All shared packages and macros belong here.
%%
%% ximera.cls already loads: amsmath, amssymb, amsthm, cancel, enumitem, graphicx,
%% hyperref, listings, pgfplots, tikz, url, xcolor. No need to re-load those.
, so this file is
%% read twice. That was harmless while it held only \usepackage lines (LaTeX skips a
%% package it has already loaded) but a second \newcommand is a hard error.
%%
%%   \blank[width]        fill-in-the-blank rule (default 2.5cm)
%%   \showwork            "show and explain your work" reminder
%%   \showworkline        ... plus which schedule line was used
%%   \showworkyear        ... plus which year's schedule was used
%%   \schedhead{...}      centred bold caption above a tiered-rate table
%%   \samesquares[scale]{left}{right}   two equal-area squares, labelled
%%   \samecubes[scale]{left}{right}     two equal-volume cubes, labelled
%% ---------------------------------------------------------------

\providecommand{\blank}[1][2.5cm]{\underline{\hspace{#1}}}
\providecommand{\showwork}{\textbf{REMEMBER TO SHOW AND EXPLAIN YOUR WORK.}}
\providecommand{\showworkline}{\textbf{REMEMBER TO SHOW AND EXPLAIN YOUR WORK.
  Explanation should include how and why you know which line of the
  schedule was used to calculate this amount.}}
\providecommand{\showworkyear}{\begin{sloppypar}\textbf{REMEMBER TO SHOW AND
  EXPLAIN YOUR WORK. Explanation should include how and why you know which
  line of the year's tax schedule was used to calculate this tax
  amount.}\end{sloppypar}}
\providecommand{\schedhead}[1]{\begin{center}\textbf{#1}\end{center}}

\providecommand{\samesquares}[3][1]{%
\begin{center}
{\footnotesize These squares are the \textbf{same size}.}

\vspace{1.5mm}
\begin{tikzpicture}[scale=#1,every node/.style={scale=#1}]
  \draw (0,0) rectangle (2.4,2.4);
  \node[anchor=north west] at (0.15,2.25) {Area:};
  \node[anchor=east]  at (-0.15,1.2) {#2};
  \node[anchor=north] at (1.2,-0.15) {#2};
  \begin{scope}[xshift=4.6cm]
    \draw (0,0) rectangle (2.4,2.4);
    \node[anchor=north west] at (0.15,2.25) {Area:};
    \node[anchor=east]  at (-0.15,1.2) {#3};
    \node[anchor=north] at (1.2,-0.15) {#3};
  \end{scope}
\end{tikzpicture}
\end{center}}

\providecommand{\samecubes}[3][1]{%
\begin{center}
{\footnotesize These cubes are the \textbf{same size}.}

\vspace{1.5mm}
\begin{tikzpicture}[scale=#1,every node/.style={scale=#1}]
  \foreach \dx in {0,5.2} {
    \begin{scope}[xshift=\dx cm]
      \draw (0,0) rectangle (2.0,2.0);
      \draw (0,2.0) -- (0.65,2.6) -- (2.65,2.6) -- (2.0,2.0);
      \fill[black!12] (2.0,0) -- (2.65,0.6) -- (2.65,2.6) -- (2.0,2.0) -- cycle;
      \draw (2.0,0) -- (2.65,0.6) -- (2.65,2.6) -- (2.0,2.0);
      \node[anchor=north west] at (0.1,1.9) {Volume:};
    \end{scope}
  }
  \node[anchor=east]  at (-0.15,1.0) {#2};
  \node[anchor=north] at (1.0,-0.15) {#2};
  \node[anchor=west]  at (2.25,0.4) {#2};
  \node[anchor=east]  at (5.05,1.0) {#3};
  \node[anchor=north] at (6.2,-0.15) {#3};
  \node[anchor=west]  at (7.45,0.4) {#3};
\end{tikzpicture}
\end{center}}

%% NOTE: do NOT add \usepackage to individual activity .tex files.
%% When a xourse inlines an activity it gobbles \documentclass and \input, but a bare
%% \usepackage still executes -- after \begin{document} -- and errors out.
%% All shared packages and macros belong here.
%%
%% ximera.cls already loads: amsmath, amssymb, amsthm, cancel, enumitem, graphicx,
%% hyperref, listings, pgfplots, tikz, url, xcolor. No need to re-load those.

\title{In-Class: Exponentials, Part 2}
\begin{document}
\begin{abstract}
Converting between the two forms of an exponential model, $y = a \cdot e^{bx}$ and
$y = a \cdot R^{x}$, using the class dice-rolling data --- and then building a
theoretical model from probability alone and comparing the two.
\end{abstract}
\maketitle

\section*{Some exponential rules}

The rule we need throughout is
\[
\left(a^{b}\right)^{c} = a^{bc}.
\]

\begin{example}
Take $a = 2$, $b = 2$, $c = 3$. Then $\left(2^{2}\right)^{3} = 2^{2 \cdot 3} = 2^{6}$.
Checking the left side directly,
\[
\left(2^{2}\right)^{3} = 2^{2} \cdot 2^{2} \cdot 2^{2}
= (2 \cdot 2)(2 \cdot 2)(2 \cdot 2) = 2^{6},
\]
so both sides agree.
\end{example}

\begin{example}
Now take $a = e$, $b = 0.3$, $c = x$, where $e \approx 2.718$ is Euler's number.
Using the same rule,
\[
e^{0.3 \cdot x} = \left(e^{0.3}\right)^{x} \approx (1.3499)^{x}.
\]
In general, given anything of the form $e^{bx}$, we can rewrite it as
$\left(e^{b}\right)^{x}$.
\end{example}

Last week all of our exponential models were written as $y = a \cdot R^{x}$. But
exponential models also come in the form $y = a \cdot e^{bx}$ --- and that second
form is actually what most programs produce. The ``$e$'' is just a letter for a
special irrational number: like $\pi = 3.14159\ldots$, we have
$e = 2.71828\ldots$. Your scientific calculator has a button for it.

\begin{problem}
Suppose we have the model $y = 671 \cdot e^{0.3x}$. Rewrite it in $y = a \cdot R^{x}$
form. Calculate the $\left(e^{b}\right)$-value rounded to \textbf{4 decimal places}.

\[
R = e^{0.3} = \answer[tolerance=0.0001]{1.3499},
\qquad
y = \answer{671} \cdot (1.3499)^{x}
\]

\begin{solution}
$e^{0.3 \cdot x} = \left(e^{0.3}\right)^{x} = (1.3499)^{x}$, so
$y = 671 \cdot e^{0.3x}$ can be written as $y = 671 \cdot (1.3499)^{x}$.
\end{solution}
\end{problem}

\section*{The dice data}

Models of the form $y = a \cdot e^{bx}$ are what Google Sheets produces for our dice
data. Open the class dice data file you completed in your pre-class work; we will
have it produce an exponential model for each of the $n$-sided dice datasets.

\begin{quote}
\textbf{Getting an exponential trendline in Google Sheets.} Select the data you want
to graph, click the ``Insert'' menu, choose ``Charts,'' and a chart appears with an
editing box on the right. At the top of that box is a drop-down to select the chart
type --- scroll down and choose the scatter plot. You can adjust the horizontal and
vertical ranges under ``Customize.''

To get the trendline: click one of the data points, and the editing box will show a
checkbox for ``trendline.'' Checking it produces a trendline and more options. It
defaults to linear, so change the ``type'' to ``exponential.''

To get the equation: in the options that open after checking ``trendline,'' there is
a drop-down called ``Label.'' Select ``Use Equation'' to display the equation on the
chart.
\end{quote}

For now let's focus on the 6-sided dice data.

\begin{center}
\begin{tabular}{|l|l|l|l|}
\hline
Dice & Model as $y = a \cdot e^{bx}$ & Model as $y = a \cdot R^{x}$ & Actual starting \# \\
\hline
6-sided & $y = 336 e^{-0.181x}$ & $y = 336 \cdot (0.834)^{x}$ & 325 \\
\hline
\end{tabular}
\end{center}

\[
e^{-0.181 \cdot x} = \left(e^{-0.181}\right)^{x} \approx (0.834)^{x}
\]

\begin{problem}
What does the $x$-value in our model represent?
\begin{freeResponse}
\end{freeResponse}
\begin{solution}
$x$ is the number of \textbf{rolls} that have occurred.
\end{solution}
\end{problem}

\begin{problem}
What does the $y$-value in our model represent?
\begin{freeResponse}
\end{freeResponse}
\begin{solution}
$y$ is the number of dice remaining after $x$ rolls.
\end{solution}
\end{problem}

\begin{problem}
The starting number of dice corresponds to $x = \answer{0}$. Use this and the model
to determine the model's estimate for the starting number of dice.
\[
\answer{336} \text{ dice}
\]

\begin{solution}
The starting number of dice is when no rolls have occurred, which is $x = 0$ since
$x$ counts rolls.

Given a model $y = a \cdot R^{x}$, when $x = 0$ we get
$y = a \cdot R^{0} = a \cdot 1 = a$. So $a$ is always the model's value at $x = 0$.
According to the model, the initial number of dice is $336$.
\end{solution}
\end{problem}

\begin{problem}
Why does the model's estimate not match our actual starting number of dice? (We
actually started with 325; the model says 336.)

\begin{freeResponse}
\end{freeResponse}

\begin{solution}
The model comes from fitting an exponential curve to the data, and the data is not
perfectly exponential --- we are only approximating the situation with an exponential
model. The fit trades off accuracy at the start against accuracy elsewhere, so its
$a$-value need not land exactly on the true starting count.
\end{solution}
\end{problem}

\begin{problem}
Plug $x = 8$ into the model $y = 336 e^{-0.181x}$. Explain what the $x$- and
$y$-values mean in this context.

$y \approx \answer[tolerance=1]{79}$ dice.

\begin{freeResponse}
\end{freeResponse}

\begin{solution}
When $x = 8$, $y = 336 \cdot e^{-0.181 \cdot 8} \approx 78.97$. So the model predicts
there will be about 79 dice left after 8 rolls.
\end{solution}
\end{problem}

\begin{problem}
Using $e^{bx} = \left(e^{b}\right)^{x}$ and the numbers from the model, calculate the
$\left(e^{b}\right)$-value, rounded to 4 decimal places, and write the model in the
form $y = a \cdot R^{x}$.
\[
R = e^{-0.181} = \answer[tolerance=0.0001]{0.8344},
\qquad
y = \answer{336} \cdot (0.834)^{x}
\]
\end{problem}

\begin{problem}
Now plug $x = 8$ into the $y = a \cdot R^{x}$ model.
\[
y \approx \answer[tolerance=1]{79} \text{ dice}
\]

Did you get the same value as with the $y = a \cdot e^{bx}$ model? Should you have?
Explain.

\begin{freeResponse}
\end{freeResponse}

\begin{solution}
$336 \cdot (0.834)^{8} = 78.644$, versus $78.97$ from the $e^{bx}$ form. Rounded to
whole dice both give 79.

They should agree, because the two forms are algebraically the same function --- we
only rewrote $e^{-0.181}$ as $0.834$. The small difference comes entirely from
rounding $e^{-0.181} = 0.8344\ldots$ down to $0.834$.
\end{solution}
\end{problem}

\begin{problem}
What should the $R$-value in the $y = a \cdot R^{x}$ model be representing in this
context?

\begin{freeResponse}
\end{freeResponse}

\begin{solution}
The $R$-value tells us about the percent change in $y$ for an increase of 1 in $x$.
For example, in the athlete's salary model $y = 8.3(1.048)^{x}$, $R = 1.048$ told us
the salary increases by $4.8\%$ each year.

Here the model is $y = 336 \cdot (0.834)^{x}$, so $R = 0.834$ --- note it is
\emph{less} than 1. This means $83.4\%$ of the dice remain after each roll, on
average. Equivalently, since $1 - 0.834 = 0.166$, we can write
\[
y = 336(1 - 0.166)^{x},
\]
which says we remove $16.6\%$ of the dice on each roll, on average.
\end{solution}
\end{problem}

\subsection*{The other dice}

Now let's look at the data for the other dice.

\begin{problem}
Fill in the $R = e^{b}$ column for the 4-, 8-, and 12-sided dice. The 6-sided row is
completed for you as a worked example.

\begin{center}
\begin{tabular}{|l|l|l|l|l|}
\hline
Dice & Model as $y = a \cdot e^{bx}$ & Actual starting \# & $R = e^{b}$ & Model as $y = a \cdot R^{x}$ \\
\hline
6-sided & $y = 336 e^{-0.181x}$ & 325 & $0.834$ & $y = 336 \cdot (0.834)^{x}$ \\
\hline
4-sided & $y = 264 e^{-0.315x}$ & 300 & $\answer[tolerance=0.005]{0.73}$ & $y = 264 \cdot (0.73)^{x}$ \\
\hline
8-sided & $y = 265 e^{-0.152x}$ & 260 & $\answer[tolerance=0.005]{0.86}$ & $y = 265 \cdot (0.86)^{x}$ \\
\hline
12-sided & $y = 161 e^{-0.0899x}$ & 165 & $\answer[tolerance=0.005]{0.914}$ & $y = 161 \cdot (0.914)^{x}$ \\
\hline
\end{tabular}
\end{center}

\begin{solution}
Each $R$ is just $e$ raised to that model's $b$:
\begin{align*}
\text{4-sided:} \quad & R = e^{-0.315} = 0.7298\ldots \approx 0.73 \\
\text{8-sided:} \quad & R = e^{-0.152} = 0.8590\ldots \approx 0.86 \\
\text{12-sided:} \quad & R = e^{-0.0899} = 0.9140\ldots \approx 0.914
\end{align*}
Checking the worked row the same way, $e^{-0.181} = 0.8344\ldots \approx 0.834$, which
matches.

Notice the ordering: the 4-sided dice have the smallest $R$ and the 12-sided the
largest. That is what you should expect --- a 4-sided die has a $1/4$ chance of showing
a ``1'' and being removed each roll, so those piles shrink fastest, while a 12-sided die
only has a $1/12$ chance, so more of it survives each round.
\end{solution}
\end{problem}

\begin{problem}
What does the $x$-value represent for all of the models?
\begin{freeResponse}
\end{freeResponse}
\begin{solution}
$x$ is the number of \textbf{rolls} that have occurred.
\end{solution}
\end{problem}

\begin{problem}
What does the $y$-value represent for all of the models?
\begin{freeResponse}
\end{freeResponse}
\begin{solution}
$y$ is the number of dice that remains after $x$ rolls.
\end{solution}
\end{problem}

\begin{problem}
What does the $a$-value represent in all of the models?
\begin{freeResponse}
\end{freeResponse}
\begin{solution}
In $y = a \cdot R^{x}$, the value $a$ is the initial number of dice, as predicted by
the model.
\end{solution}
\end{problem}

\begin{problem}
What does the $R$-value (which equals $e^{b}$) represent in all of the models?
\begin{freeResponse}
\end{freeResponse}
\begin{solution}
$R$ is the proportion of dice that remain after each roll, and $1 - R$ is the
proportion removed after each roll.
\end{solution}
\end{problem}

\section*{Theoretical exponential models}

Excel and Google Sheets give models in the form $y = a \cdot e^{bx}$, and we rewrote
them as $y = a \cdot R^{x}$ to see the connection to the context. But without rolling
any dice at all, and without any special program, we could have built a very decent
model ourselves. Let's see how.

Do a thought experiment with 6-sided dice, following the same procedure as before: we
roll the dice and remove any that come up ``1.''

\begin{problem}
Theoretically, what fraction of the dice should roll a ``1''?
\[
\answer{1/6}
\]
\begin{solution}
Assuming each side is equally likely, there are six possibilities, so there is a 1 in
6 chance of rolling a 1. We expect about 1 out of every 6 dice to roll a 1.
\end{solution}
\end{problem}

\begin{problem}
Theoretically, what fraction of the dice should \emph{not} roll a ``1''?
\[
\answer{5/6}
\]
\begin{solution}
Since $1/6$ of the dice should roll a 1, we expect $5/6$ of the dice not to.
\end{solution}
\end{problem}

\begin{problem}
Suppose we start with the same number of 6-sided dice as your group's actual starting
number: $325$.

\begin{prompt}
\textbf{(a)} In a theoretical model $y = a \cdot R^{x}$, which value does the
starting number of dice go in for, and why?
\begin{freeResponse}
\end{freeResponse}
\begin{solution}
It is the $a$-value, since $a$ is the model's output when $x = 0$ --- before any
rolls have happened.
\end{solution}
\end{prompt}

\begin{prompt}
\textbf{(b)} Which value should go in as the $R$-value, and why?
\begin{freeResponse}
\end{freeResponse}
\begin{solution}
$R = 5/6$, the proportion of dice we expect to survive each roll.
\end{solution}
\end{prompt}

\begin{prompt}
\textbf{(c)} Put these together to create the theoretical model for our 6-sided dice.
\[
y = \answer{325} \cdot \left(\answer{5/6}\right)^{x}
\]
\end{prompt}
\end{problem}

\begin{problem}
Look back at the model the program produced for the 6-sided dice in its
$y = a \cdot R^{x}$ form. Its $R$-value should be close to $0.8333\ldots$ --- why
would we hope and expect that?

\begin{freeResponse}
\end{freeResponse}

\begin{solution}
Because $5/6 = 0.8333\ldots$ is the theoretical survival proportion. The program's
$R$ came from the actual rolls, and if the dice are fair, the observed proportion
surviving each roll should hover near the theoretical one. The program got $0.834$,
which is very close indeed.
\end{solution}
\end{problem}

\begin{problem}
Compare the two models using the chart below. Fill in the ``actual'' column from your
group's data, and round the model values to the nearest whole number.

\begin{center}
\begin{tabular}{|c|l|l|l|}
\hline
Roll & Actual \# remaining & Google Sheets model $336(0.834)^{x}$ & Theoretical model $325(5/6)^{x}$ \\
\hline
0 & & $\answer{336}$ & $\answer{325}$ \\
\hline
1 & & $\answer[tolerance=1]{280}$ & $\answer[tolerance=1]{271}$ \\
\hline
2 & & $\answer[tolerance=1]{234}$ & $\answer[tolerance=1]{226}$ \\
\hline
3 & & $\answer[tolerance=1]{195}$ & $\answer[tolerance=1]{188}$ \\
\hline
4 & & $\answer[tolerance=1]{163}$ & $\answer[tolerance=1]{157}$ \\
\hline
\end{tabular}
\end{center}
\end{problem}

\begin{problem}
How do these values compare? Which is better, the Google Sheets model or the
theoretical model? Which would we \emph{expect} to be better, and why?

\begin{freeResponse}
\end{freeResponse}
\end{problem}

\end{document}
